\documentclass[11pt]{article}

\usepackage[utf8]{inputenc}
\usepackage[margin=1in]{geometry}
\usepackage{amsmath, amssymb, amsthm}
\usepackage{float}
\usepackage{algorithm}
\usepackage[noend]{algpseudocode}

\newtheorem{theorem}{Theorem}
\newtheorem{lemma}[theorem]{Lemma}

\newcommand{\R}{\mathbb{R}}
\newcommand{\norm}[1]{\left\| #1 \right\|}
\newcommand{\inner}[2]{\left\langle #1, #2 \right\rangle}
\DeclareMathOperator*{\argmin}{arg\,min}

\title{Different methods for gradient descent stepsize choice}
\author{}
\date{}

\begin{document}

\maketitle

Let us assume that we are given a function $f\colon \R^d \to \R$ which is $L$-smooth,
$\mu$-strongly convex, and twice continuously differentiable, i.e., $f \in C^2(\R^d)$.
As a reminder, Gradient Descent (GD) is an iterative method for finding the minimum
of a function, and it can be described as follows.

\begin{algorithm}[H]
\caption{GD}
\begin{algorithmic}[1]
\State \textbf{Parameters:} stepsize $\gamma > 0$, starting point $x_0 \in \R^d$
\For{$k = 0, 1, 2, \dots$}
  \State Set $g_k = \nabla f(x_k)$
  \State $x_{k+1} = x_k - \gamma \nabla f(x_k)$
\EndFor
\end{algorithmic}
\end{algorithm}

Our aim is to choose the stepsize $\gamma$ so that GD converges ``faster''.
The connection between the stepsize and the convergence speed can be described by the following lemma.

\begin{lemma}\label{lem:one-step}
Assume $f$ is $L$-smooth and $\mu$-strongly convex. Then the iterates of GD satisfy
the following inequality for all $\gamma \geq 0$ and all $k \geq 0$:
\begin{equation}\label{eq:lemma}
\norm{x_{k+1}-x^*}^2 \leq (1-\gamma\mu)\norm{x_k-x^*}^2 - 2\gamma (1 - \gamma L)\, D_f(x_k, x^*),
\end{equation}
where $x^* \overset{\text{def}}{=} \argmin\limits_{x\in\R^d} f(x)$ and $D_f(x,y)$ is the Bregman
divergence, defined as follows
\[
D_f(x, y) \overset{\text{def}}{=} f(x) - f(y) - \inner{\nabla f(y)}{x - y}
\qquad \forall x, y \in \R^d.
\]
\end{lemma}

Note that $D_f(x,y) \geq 0$ for every convex function $f$, and hence for every
$\mu$-strongly convex function.

\begin{proof}[Proof of Lemma~\ref{lem:one-step}]
\begin{align*}
\norm{x_{k+1}-x^*}^2
&= \norm{\bigl(x_k - \gamma\nabla f(x_k)\bigr)-\bigl(x^* - \gamma\nabla f(x^*)\bigr)}^2 \\
&= \norm{(x_k - x^*) - \gamma\bigl(\nabla f(x_k) -\nabla f(x^*)\bigr)}^2 \\
&= \norm{x_k - x^*}^2 + \gamma^2\norm{\nabla f(x_k) -\nabla f(x^*)}^2
   - 2\gamma\inner{x_k - x^*}{\nabla f(x_k) -\nabla f(x^*)}.
\end{align*}
From strong convexity we obtain
\[
\inner{x_k - x^*}{\nabla f(x_k) -\nabla f(x^*)} \geq D_f(x_k,x^*) + \frac{\mu}{2}\norm{x_k - x^*}^2.
\]
Also, from $L$-smoothness we have
\[
\norm{\nabla f(x_k) -\nabla f(x^*)}^2 \leq 2L\, D_f(x_k,x^*),
\]
and, as a result,
\[
\norm{x_{k+1}-x^*}^2 \leq (1-\gamma\mu)\norm{x_k-x^*}^2 - 2\gamma (1 - \gamma L)\, D_f(x_k, x^*),
\]
which finalizes the proof.
\end{proof}

As mentioned above, $D_f(x_k, x^*) \geq 0$. Hence, if we take $0 \leq \gamma \leq 1/L$,
then $1 - \gamma L \geq 0$ and the last term in \eqref{eq:lemma} is non-positive.
Dropping it and applying \eqref{eq:lemma} $k$ times, we obtain the following theorem.

\begin{theorem}[Convergence of GD]\label{thm:gd}
Assume $f$ is $L$-smooth and $\mu$-strongly convex. Then the iterates of GD
with stepsize $0 \leq \gamma \leq 1/L$ satisfy
\begin{equation}\label{eq:gd-rate}
\norm{x_k-x^*}^2 \leq (1-\gamma\mu)^k\norm{x_0-x^*}^2.
\end{equation}
\end{theorem}

To make this result more concrete,
we may ask how many GD iterations it takes for the error to reach
a certain target level $\varepsilon$. To answer this, one can start from inequality
\eqref{eq:gd-rate} and show that with $\gamma = 1/L$, for every $k$ that satisfies
\[
k \geq \frac{L}{\mu}\log\frac{1}{\varepsilon}
\]
we have
\[
\norm{x_k-x^*}^2 \leq \varepsilon\norm{x_0-x^*}^2.
\]

What we can conclude is that if the function is $L$-smooth and
$\mu$-strongly convex, then we can choose some $0 \leq \gamma \leq 1/L$ in \textbf{advance}
for all iterations of GD, and we can guarantee that after a certain number of steps
the squared distance between the current point and the minimizer is reduced by a factor
of $\frac{1}{\varepsilon}$ compared to the initial point.

But what if we are allowed to change the stepsize during the algorithm,
that is, to pick a new stepsize $\gamma_k$ at each step $k$?
If we are willing to do this ``extra'' work, and of course if we are able to do it,
can we find stepsizes $\gamma_k$ that are better than a constant one? Or at least not worse, in the sense of the $\varepsilon$-guarantee
from the previous part?

It turns out that the answer is yes, and one such method is called the Polyak stepsize.

Consider a variant of GD with Polyak stepsize, i.e., with the adaptive (i.e.,
iteration-dependent) stepsize rule
\[
\gamma_k = \frac{2D_f(x_k,x^*)}{\norm{\nabla f(x_k) - \nabla f(x^*)}^2},
\]
where we assume that $\nabla f(x_k) \neq \nabla f(x^*)$. This is a very natural assumption:
since $\nabla f(x^*) = 0$ and the minimizer of a strongly convex function is unique, so
$\nabla f(x_k) = \nabla f(x^*)$ would mean that $x_k = x^*$, i.e., we have already found the minimum.

\begin{enumerate}
  \item[(i)] From $L$-smoothness it can be obtained that $\gamma_k \geq 1/L$. That is,
  the Polyak stepsize is always at least as large as the largest constant stepsize
  $\gamma = 1/L$ allowed by Theorem~\ref{thm:gd}.
  \item[(ii)] From $\mu$-strong convexity it can be obtained that $\gamma_k \leq 1/\mu$.
  \item[(iii)] Finally, it can be shown that after one step of GD with Polyak stepsize we have
  \[
  \norm{x_{k+1}-x^*}^2 \leq (1-\gamma_k\mu)\norm{x_k-x^*}^2.
  \]
\end{enumerate}

Indeed, let us start from the GD step $x_{k+1} = x_k - \gamma_k\nabla f(x_k)$.
Since $\nabla f(x^*) = 0$, we can subtract the zero term $\gamma_k\nabla f(x^*)$
without changing anything, and then group the terms:
\begin{align*}
\norm{x_{k+1}-x^*}^2
&= \norm{\bigl(x_k - \gamma_k\nabla f(x_k)\bigr)-\bigl(x^* - \gamma_k\nabla f(x^*)\bigr)}^2
 = \norm{(x_k - x^*) - \gamma_k\bigl(\nabla f(x_k) -\nabla f(x^*)\bigr)}^2.
\end{align*}
Expanding the square and taking $\gamma_k$ out of the last two terms, we obtain
\begin{align*}
\norm{x_{k+1}-x^*}^2
&= \norm{x_k - x^*}^2 + \gamma_k\Bigl(\gamma_k\norm{\nabla f(x_k) -\nabla f(x^*)}^2
   - 2\inner{x_k - x^*}{\nabla f(x_k) -\nabla f(x^*)}\Bigr).
\end{align*}
Now we bound the inner product from below. By $\mu$-strong convexity,
\[
\inner{x_k - x^*}{\nabla f(x_k) -\nabla f(x^*)} \geq D_f(x_k,x^*) + \frac{\mu}{2}\norm{x_k - x^*}^2.
\]
The inner product enters with a minus sign and $\gamma_k \geq 0$, so replacing it by this
lower bound can only increase the right-hand side. Moreover, in front of the squared gradient
norm we substitute the definition of $\gamma_k$. This gives
\begin{align*}
\norm{x_{k+1}-x^*}^2 \leq \norm{x_k - x^*}^2 + \gamma_k\Biggl(\frac{2D_f(x_k,x^*)}{\norm{\nabla f(x_k) -\nabla f(x^*)}^2}
   \cdot\norm{\nabla f(x_k) -\nabla f(x^*)}^2  - 2D_f(x_k,x^*) - \mu\norm{x_k - x^*}^2\Biggr).
\end{align*}
Here the squared gradient norms cancel, so the first term in the brackets is just
$2D_f(x_k,x^*)$, and it cancels with $-2D_f(x_k,x^*)$. Only the term
$-\mu\norm{x_k - x^*}^2$ remains, and therefore
\[
\norm{x_{k+1}-x^*}^2 \leq \norm{x_k - x^*}^2 - \gamma_k\mu\norm{x_k - x^*}^2
= (1-\gamma_k\mu)\norm{x_k - x^*}^2,
\]
which is exactly the inequality in (iii).

\end{document}
